\section{Texture, pores and the retention curve}
The texture triangle places a soil by its three fractions; the retention curve shows the water held at each suction. A \kw{pore domain} is a class of pore sizes that has a function in the soil. The \emph{macropores} (radius above about 30 $\mu$m, drained at less than 100 cm of suction, pF below 2) carry the fast flow and the air; the \emph{mesopores} (30 to 0.2 $\mu$m, pF 2 to 4.2) hold the water that the plants use; the \emph{micropores} (below 0.2 $\mu$m, pF above 4.2) hold water that the plants cannot take. The sizes follow from the capillary law $r=0.149/h$ ($r$ and $h$ in cm). The volume of each domain is the difference of the water contents at its limits, so that the pore-domain figure is a direct reading of the retention curve.

\section{Correlation, importance and the random forest}
A \kw{correlation matrix} shows, for each pair of variables, the rank correlation $\rho$; it is a quick map of which quantities move together, but it describes pairs only. The \kw{random forest} is a set of decision trees, each fitted to a random sample of the runs and a random subset of the variables; the average of the trees predicts the outcome, and the \emph{importance} of a variable is the total reduction of the error of the trees that comes from splits on it. The share of the variance explained is checked on the runs that each tree did not see (the out-of-bag $R^2$). An importance is a ranking of the drivers, not a measure of their effect, and correlated drivers share the importance among them.

\section{Principal components}
When there are many variables that move together, the \kw{principal component analysis} (PCA) finds new axes, the principal components, which are combinations of the variables, ordered by the variance they explain. The data are first standardised ($z=(x-\mu)/\sigma$) so that the units do not matter. The first axis is the direction of the largest variance; the second is the largest in the direction at right angles to the first. The \kw{biplot} shows the runs as points and the variables as arrows: arrows that point the same way are positively correlated, arrows at right angles are unrelated, and opposite arrows are negatively correlated. The ratio $\lambda_j/\sum\lambda$ is the share of the variance of the component.

\section{Parallel coordinates}
In a \kw{parallel-coordinates} plot each variable is a vertical axis, scaled from its minimum to its maximum, and each run is a line joining its values. The lines of the best runs, chosen by a criterion, show what values of each variable they share. It allows many variables to be read at once, but the order of the axes changes what is seen.

\section{Variance decomposition}
The variance of an outcome can be divided among the factors. For a factor with levels $\ell$, the sum of squares $\mathrm{SS}_k=\sum_\ell n_\ell(\bar y_\ell-\bar y)^2$ measures how much the level means differ, and the share $\mathrm{SS}_k/\mathrm{SS}_{tot}$ is the fraction of the variance that the factor explains. In a balanced design the shares add up to a total of at most one; what remains is the \emph{unexplained} part, which includes the interactions. A large unexplained share means the factors act in combination, not alone, and the interaction plots of the previous chapter are the next step.

\section{Validation by re-running}
The daily figures need the daily output, which the batch of runs did not keep, so 36 setups were run again with a later version of the program. A re-run is a test of reproducibility: the correlation between the re-run and the original, and the median relative difference, say how much the new version changes the results. A median difference below 0.3 per cent, with a worst case below 2 per cent, means that the daily figures can be used as a picture of the batch.
